% !TEX root = ../main.tex
\chapter{等式・書き換え・正規化}
\label{chap:ch18_equations_rewriting_normalization}
\BookIndex{かきかえ}{書き換え}
\BookIndex{せいきか}{正規化}
\BookIndex{calc}{\texttt{calc}}

\begin{chapterabstract}
本章では、Lean 4 における等式証明を、仕様変更とリファクタリングの安全性を確認する道具として扱います。
中心になる道具は、\texttt{rw}、\texttt{simp}、\texttt{calc} です。
\texttt{rw} は既知の等式に基づく置換、\texttt{simp} は自明な整理による正規化、\texttt{calc} は人間がレビューしやすい等式変形のログとして読みます。
本章のサンプルでは、価格計算ロジックを小さくモデル化し、関数分割、補助関数の導入、手数料計算の括弧変更が意味を変えないことを Lean に確認させます。
\end{chapterabstract}

\begin{learninggoals}
\item 等式証明を「変更前後で同じ結果になること」の確認として読めます。
\item \texttt{rw}、\texttt{simp}、\texttt{calc} の役割を区別できます。
\item 価格計算のような小さな業務ロジックを、等式として整理できます。
\item 正規化を、レビューと自動証明の両方に有効な設計方針として使えます。
\end{learninggoals}

\section{問題設定}

実務で頻繁に起きる変更の多くには、外から見た意味を変えないことが求められます。
たとえば、関数の分割、補助関数の導入、処理の共通化、括弧の位置の変更、不要な初期値処理の削除です。
これらは、「同じ結果を返す別の書き方」への変更として捉えられます。

Lean では、この主張を等式として書きます。

\begin{center}
\begin{tabular}{lll}
\toprule
実務上の言い方 & Lean での形 & 本章で使う道具 \\
\midrule
変更前後で結果が同じ & \texttt{old x = new x} & \texttt{rfl}, \texttt{rw}, \texttt{simp} \\
既知の仕様で置き換える & \texttt{h : a = b} を使う & \texttt{rw [h]} \\
自明な整理を自動化する & \texttt{x + 0} を \texttt{x} にする & \texttt{simp} \\
レビュー可能に段階を書く & 等式を一行ずつつなぐ & \texttt{calc} \\
\bottomrule
\end{tabular}
\end{center}

図\ref{fig:ch18-rewriting-flow}では、変更前の式を検査済みの等式で段階的に変形し、変更後と同じ形へ到達します。
各矢印に使った補題を明示すれば、証明を変更理由のログとして読めます。

\FigurePlaceholder{fig-ch18-rewriting-flow.png}{0.92}{等式による書き換え}{fig:ch18-rewriting-flow}

\section{仕様証明の多くは等式になる}

仕様には、不変条件、事前条件、事後条件、順序関係、到達可能性、エラーケースの網羅など、さまざまな形があります。
その中でも、リファクタリングやデータ変換で最初に現れるのは等式です。

たとえば、次のような主張はすべて等式として表せます。

\begin{itemize}
\item 料金計算を分割しても、合計金額は変わりません。
\item 何もしない正規化処理を二回通しても、一回通した結果と変わりません。
\item 手数料を一つずつ足しても、先にまとめてから足しても同じです。
\item 古い関数を新しい関数に置き換えても、すべての入力で同じ出力になります。
\end{itemize}

この章では、本番コード全体ではなく、計算ロジックの核だけを Lean に移します。
小数、丸め、通貨単位、外部設定などを省いた範囲を明示し、そのモデル内での書き換えを検証します。

\section{最初のモデル：価格計算を小さく切り出す}

最初に、割引後の小計に税を足すだけの小さな価格計算を考えます。
ここでは \texttt{Nat} を使います。
\texttt{Nat} は自然数であり、負の値を扱いません。
実務の金額型としては単純ですが、等式証明の形を学ぶには十分です。

\begin{listing}[htbp]
\begin{minted}{lean4}
namespace Chapter18

-- 割引後の小計。Nat の引き算は 0 で止まる。
def subtotal (base discount : Nat) : Nat :=
  base - discount

-- 税や手数料を加える単純な処理。
def addTax (amount tax : Nat) : Nat :=
  amount + tax

def legacyPrice (base discount tax : Nat) : Nat :=
  addTax (subtotal base discount) tax

def refactoredPrice (base discount tax : Nat) : Nat :=
  (base - discount) + tax

#eval legacyPrice 1000 100 80  -- => 980
#eval refactoredPrice 1000 100 80  -- => 980

theorem legacy_eq_refactored
    (base discount tax : Nat) :
    legacyPrice base discount tax =
      refactoredPrice base discount tax := by
  rfl
\end{minted}
\caption{\texttt{legacyPrice\_eq\_refactored}}
\label{lst:ch18-price-model}
\end{listing}

\texttt{rfl} は、反射律を使うタクティクです。
直感的には、「両辺を定義に従って展開すると同じ形になる」と Lean が確認できる場合に使います。
上の例では、\texttt{legacyPrice} と \texttt{refactoredPrice} は見た目が異なりますが、定義を展開すると同じ式になります。

\begin{warningbox}
\texttt{rfl} で証明できるからといって、その仕様が実務上十分だとは限りません。
このモデルでは、丸め、税率、負の金額、通貨単位、上限と下限を扱っていません。
本章で証明しているのは、「この小さなモデルの範囲では変更前後が同じ」という主張です。
\end{warningbox}

\section{\texttt{rw}：仕様に基づく置換}

\texttt{rw} は rewrite の略です。
すでに持っている等式を使って、証明ゴールの中にある式を置き換えます。

\begin{definitionbox}
等式 \texttt{h : a = b} があるとき、\texttt{rw [h]} はゴール中の \texttt{a} を \texttt{b} に置き換えます。
逆向きに使う場合は、\texttt{rw [← h]} と書きます。
\end{definitionbox}

この見方は、仕様レビューにそのまま対応します。
補題 \texttt{addTax\_zero} は、「税額が 0 なら、税を足しても金額は変わらない」という小さな仕様です。
その仕様を使って、\texttt{normalizedPrice} を \texttt{subtotal} に書き換えます。

\begin{listing}[htbp]
\begin{minted}{lean4}
-- 0 を足す税計算は、元の金額と同じ。
theorem addTax_zero (n : Nat) :
    addTax n 0 = n := by
  unfold addTax
  rw [Nat.add_zero]

-- 既存コードには、不要な 0 税計算が残っているとする。
def normalizedPrice (base discount : Nat) : Nat :=
  addTax (subtotal base discount) 0

-- その不要な処理を仕様に基づいて取り除く。
theorem normalizedPrice_eq_subtotal
    (base discount : Nat) :
    normalizedPrice base discount =
      subtotal base discount := by
  unfold normalizedPrice
  rw [addTax_zero]
\end{minted}
\caption{\texttt{rw} による仕様ベースの置換}
\label{lst:ch18-rw}
\end{listing}

\texttt{rw} は「文字列置換」ではありません。
\texttt{rw} が使うのは、Lean が検査済みの等式です。
つまり、置換には根拠があります。
コードレビューでは、「この変更は \texttt{addTax\_zero} という仕様に基づきます」と説明できます。

\subsection{逆向きの書き換え}

より複雑な形へ戻したいこともあります。
たとえば、既存のインターフェースに合わせるため、一時的に \texttt{n} を \texttt{addTax n 0} として見る場合です。
このときは、等式を逆向きに使います。

\begin{listing}[htbp]
\begin{minted}{lean4}
theorem subtotal_eq_normalizedPrice
    (base discount : Nat) :
    subtotal base discount =
      normalizedPrice base discount := by
  rw [← normalizedPrice_eq_subtotal]
\end{minted}
\caption{\texttt{rw [← ...]}}
\label{lst:ch18-rw-backward}
\end{listing}

この例では \texttt{normalizedPrice\_eq\_subtotal} を逆向きに使い、左辺の
\texttt{subtotal base discount} を \texttt{normalizedPrice base discount} へ書き換えます。
最後は、同じ式どうしの等式になります。

\section{\texttt{simp}：自明な整理を正規化として使う}

\texttt{rw} では、どの等式で書き換えるかを人間が明示します。
一方、\texttt{simp} は Lean が知っている単純化規則をまとめて使います。
\texttt{x + 0 = x}、\texttt{0 + x = x}、\texttt{if true then a else b = a} のような明らかな整理を毎回手で書くと、証明が読みづらくなります。

\begin{intuitionbox}
\texttt{simp} は、「式を読みやすい標準形へ寄せる道具」と考えられます。
つまり、正規化のための自動整理です。
実務のリファクタリングでも、不要な wrapper、空の追加、恒等処理を消して、比較しやすい形へ寄せることがあります。
\end{intuitionbox}

\begin{listing}[htbp]
\begin{minted}{lean4}
-- この補題は simp が自動で使えるようにしておく。
@[simp] theorem addTax_zero_simp (n : Nat) :
    addTax n 0 = n := by
  unfold addTax
  simp

-- 何もしない正規化処理を、あえて関数として置く。
def normalizeCode (n : Nat) : Nat :=
  n + 0

def loadCode (n : Nat) : Nat :=
  normalizeCode (normalizeCode n)

@[simp] theorem normalizeCode_eq (n : Nat) :
    normalizeCode n = n := by
  unfold normalizeCode
  simp

theorem loadCode_normalized (n : Nat) :
    loadCode n = n := by
  simp [loadCode]
\end{minted}
\caption{\texttt{normalizeCode\_eq}}
\label{lst:ch18-simp}
\end{listing}

\texttt{@[simp]} は、証明済みの補題を \texttt{simp} の単純化規則として登録する属性であり、それ自体が証明を作るわけではありません。
上の例では、\texttt{normalizeCode} が恒等処理であることを一度証明し、二重の呼び出しも同じ規則で取り除きます。

\begin{warningbox}
\texttt{simp} に何でも登録すればよいわけではありません。
向きの悪い規則や、式を大きくする規則を登録すると、証明が読みにくくなったり、単純化が期待どおりに進まなくなったりします。
本章では、「明らかに小さくなる」「正規形に近づく」規則だけを登録します。
\end{warningbox}

\section{\texttt{calc}：レビュー可能な証明ログ}

\texttt{simp} は短く強力ですが、何が起きたかが見えにくいことがあります。
仕様レビューで重要な証明では、\texttt{calc} で途中式を残すほうが適しています。

\begin{definitionbox}
\texttt{calc} は、等式や不等式を段階的につなぐための構文です。
各行に「なぜその変形が正しいか」を書けるため、人間が読む証明ログとして使いやすくなります。
\end{definitionbox}

次の例では、二つの手数料を順に足す既存実装と、先に手数料をまとめてから足す新実装が同じであることを示します。
根拠は、足し算の結合律です。

\begin{listing}[htbp]
\begin{minted}{lean4}
def addFee (amount fee : Nat) : Nat :=
  amount + fee

def legacyWithTwoFees
    (base service shipping : Nat) : Nat :=
  addFee (addFee base service) shipping

def refactoredWithTwoFees
    (base service shipping : Nat) : Nat :=
  addFee base (service + shipping)

theorem twoFees_eq_calc
    (base service shipping : Nat) :
    legacyWithTwoFees base service shipping =
      refactoredWithTwoFees base service shipping := by
  calc
    legacyWithTwoFees base service shipping
        = (base + service) + shipping := by rfl
    _ = base + (service + shipping) := by
          rw [Nat.add_assoc]
    _ = refactoredWithTwoFees base service shipping := by rfl
\end{minted}
\caption{\texttt{twoFees\_eq\_calc}}
\label{lst:ch18-calc}
\end{listing}

この証明は、同じ定義展開と結合律を \texttt{simp} に渡して、さらに短く書けるかもしれません。
しかし、レビュー目的であれば、\texttt{calc} のほうが「どこで結合律を使ったか」を明確に示せます。

\section{正規化：比較しやすい形へ寄せる}

正規化とは、同じ意味を持つ複数の書き方を、比較しやすい代表形へ寄せることです。
Lean の \texttt{simp} は、この作業を支援します。
実務でも、比較しやすい形にしてから差分を見ると、レビューしやすくなります。

\begin{definitionbox}
本章でいう正規形とは、次のような形です。
\begin{itemize}
\item 不要な \texttt{+ 0} や恒等処理が消えています。
\item 補助関数が必要な範囲だけ展開されています。
\item 括弧の付き方が一貫しています。
\item 変更前と変更後を同じ構造で比較できます。
\end{itemize}
\end{definitionbox}

図\ref{fig:ch18-price-normalization}は、変更前後の式から不要な \texttt{+ 0} を除き、加算の括弧をそろえて共通の式へ到達する流れを示します。

\FigurePlaceholder{fig-ch18-price-normalization.png}{0.90}{価格式の正規化}{fig:ch18-price-normalization}

次の例では、既存コードが小計、二つの手数料、税額 0 の追加を段階的に行っています。
新コードでは、税額 0 を消し、二つの手数料をまとめています。
\texttt{simp} と結合律を使うと、両者を同じ正規形へ整理できます。

\begin{listing}[htbp]
\begin{minted}{lean4}
def chargeLegacy
    (base discount service shipping : Nat) : Nat :=
  addTax
    (addFee (addFee (subtotal base discount) service) shipping)
    0

def chargeNormalized
    (base discount service shipping : Nat) : Nat :=
  (base - discount) + (service + shipping)

theorem chargeLegacy_eq_chargeNormalized
    (base discount service shipping : Nat) :
    chargeLegacy base discount service shipping =
      chargeNormalized base discount service shipping := by
  unfold chargeLegacy chargeNormalized subtotal addFee addTax
  simp [Nat.add_assoc]
\end{minted}
\caption{\texttt{chargeLegacy\_eq\_chargeNormalized}}
\label{lst:ch18-price-normalization}
\end{listing}

この証明は、次のように読みます。

\begin{enumerate}
\item まず \texttt{unfold} で、業務名の付いた補助関数を計算式へ展開します。
\item 次に \texttt{simp} で、\texttt{+ 0} のような自明な部分を消します。
\item 最後に \texttt{Nat.add\_assoc} で、足し算の括弧を新実装の形へそろえます。
\end{enumerate}

ここでの \texttt{unfold} は、抽象化を一時的にほどく操作です。
通常は \texttt{subtotal} や \texttt{addFee} の名前を使うほうが読みやすくなります。
しかし、変更前後が本当に同じか比較するときには、必要な範囲だけ定義を展開します。

\section{\texorpdfstring{\texttt{rw}、\texttt{simp}、\texttt{calc} の使い分け}{rw、simp、calc の使い分け}}

三つの道具には重なる部分がありますが、目的は少し異なります。

\begin{center}
\begin{tabular}{p{0.18\linewidth}p{0.34\linewidth}p{0.36\linewidth}}
\toprule
道具 & 向いている場面 & 実務上の読み方 \\
\midrule
\texttt{rw} & 特定の補題を根拠に一箇所を書き換えます。
& この変更は、この仕様に基づく置換です。 \\
\texttt{simp} & 自明な整理をまとめて行います。
& 不要なノイズを消して正規形へ寄せます。 \\
\texttt{calc} & 途中式を残して説明する場合に使います。
& レビュー可能な変更ログを書きます。 \\
\bottomrule
\end{tabular}
\end{center}

自動化の短さと、説明の明瞭さは別の基準です。
仕様レビューで変形理由が重要なら、\texttt{calc} で途中式を残します。

\section{よくある誤解}

\subsection{誤解1：\texttt{simp} は何でも証明してくれる}

\texttt{simp} は強力ですが、万能ではありません。
基本的には、登録された単純化規則や定義展開によって式を整理する道具です。
複雑な業務仕様を自動的に理解するわけではありません。
\texttt{simp} に使わせたい業務上の事実があるなら、それを補題として先に証明する必要があります。

\subsection{誤解2：\texttt{rw} はプログラムの文字列置換である}

\texttt{rw} は、テキストエディタの置換ではありません。
Lean が型検査した等式だけを使います。
置換してよい理由が証明オブジェクトとして存在するため、単なる検索置換より厳密です。

\subsection{誤解3：短い証明ほど良い証明である}

短い証明は、保守しやすいことが多くあります。
しかし、仕様説明として途中式が必要な場合もあります。
特に設計レビューでは、\texttt{calc} による長めの証明が、将来の読者にとって最も読みやすいことがあります。

\section{実務へ持ち込むときの手順}

本章の考え方を実務のリファクタリングに使う場合、最初から大きなコードベース全体を Lean に移す必要はありません。
次の順で進めると扱いやすくなります。

\begin{enumerate}
\item 変更で守りたい性質を一文で書きます。
例：「手数料を先にまとめても合計は変わりません」。
\item その性質に関係する入力と出力だけを型として切り出します。
\item 変更前と変更後を、小さな Lean 関数として定義します。
\item 明らかな仕様を補題にします。
例：\texttt{addTax n 0 = n}。
\item \texttt{rw}、\texttt{simp}、\texttt{calc} で等式を証明します。
\item 証明した範囲と、実装側に残る前提を文書に書きます。
\end{enumerate}

\begin{warningbox}
Lean の証明モデルと本番実装は同じではありません。
特に、丸め、小数、タイムゾーン、外部設定、DB の NULL、例外処理は、単純な \texttt{Nat} モデルでは表せないことが多くあります。
証明済みと言えるのは、Lean に書いたモデルと前提の範囲に限られます。
\end{warningbox}

\section{本章のまとめ}

リファクタリングや計算式の整理は、多くの場合、「変更前後が等しい」という主張として表せます。

\texttt{rw} は、検査済みの等式を使う根拠付きの置換です。

\texttt{simp} は、不要なノイズを消して式を正規形へ寄せる道具です。

\texttt{calc} は、途中式と変形理由を残すレビュー可能な証明ログです。

正規化を意識すると、変更前後のロジックを比較しやすくなります。

Lean で証明した範囲と、実務コード側に残る前提は必ず分けて書きます。

等式の道具は、\texttt{Bool}、\texttt{Option}、\texttt{Sum} のような分岐を含む型にも適用できます。
等式だけで済む例では、\texttt{rw}、\texttt{simp}、\texttt{calc} が中心でした。
分岐がある仕様では、すべてのケースを漏れなく扱うために \texttt{cases} が必要になります。
